\documentclass{beamer}
%\documentclass[aspectratio=169]{beamer} % 16:9
\usepackage{hyperref}
\usepackage[russian]{babel}
\usepackage{fontspec}

\usepackage{latexsym,amsmath,xcolor,multicol,booktabs,calligra}
\usepackage{graphicx,pstricks,listings,stackengine}
\usepackage{array}

% Ensure Cyrillic glyph coverage under XeLaTeX.
\setmainfont{Liberation Serif}
\setsansfont{Liberation Sans}
\setmonofont{Liberation Mono}

\author{Тагир Гараев}
\title{Препроцессор C++: власть над текстом и ее пределы}
\subtitle{Архитектурный разбор + сравнение с Rust и Go}
\institute{C++ Семинар}
\date{\today}
\usepackage{CNU}

% Disable auto TOC frame on each section from CNU.sty.
\AtBeginSection[]{}
\AtBeginSubsection[]{}

\definecolor{deepblue}{rgb}{0,0,0.5}
\definecolor{deepred}{rgb}{0.6,0,0}
\definecolor{deepgreen}{rgb}{0,0.5,0}
\definecolor{halfgray}{gray}{0.55}

\lstset{
    basicstyle=\ttfamily\small,
    keywordstyle=\bfseries\color{deepblue},
    emphstyle=\ttfamily\color{deepred},
    stringstyle=\color{deepgreen},
    breaklines=true,
    columns=fullflexible,
    numbers=left,
    numberstyle=\small\color{halfgray},
    rulesepcolor=\color{red!20!green!20!blue!20},
    frame=shadowbox,
    language=C++
}

\lstdefinelanguage{Rust}{
    morekeywords={as,break,const,continue,crate,else,enum,extern,false,fn,for,if,impl,in,let,loop,match,mod,move,mut,pub,ref,return,self,Self,static,struct,super,trait,true,type,unsafe,use,where,while,macro_rules,println},
    morecomment=[l]{//},
    morecomment=[s]{/*}{*/},
    morestring=[b]",
    sensitive=true
}

\lstdefinelanguage{Go}{
    morekeywords={break,case,chan,const,continue,default,defer,else,fallback,for,func,go,if,import,interface,map,package,range,return,select,struct,switch,type,var},
    morecomment=[l]{//},
    morecomment=[s]{/*}{*/},
    morestring=[b]",
    sensitive=true
}

\begin{document}

\begin{frame}
    \titlepage
\end{frame}

\begin{frame}{План}
    \tableofcontents
\end{frame}

\section{Часть 1: Физика препроцессора}

\begin{frame}{Вводный тезис: две разные машины внутри компиляции}
    \begin{itemize}
        \item Современный C++ опирается на типы, шаблоны, concepts и AST-семантику.
        \item Но до этого запускается отдельный механизм: препроцессор (фаза 4).
        \item Это язык переписывания токенов: без типов, без областей видимости, без AST.
        \item Наша задача: понять, где его сила инженерно оправдана, а где начинается саботаж.
    \end{itemize}
\end{frame}

\begin{frame}{Конвейер трансляции: где заканчивается текст и начинается смысл}
    \begin{itemize}
        \item Стандарт задает однонаправленный конвейер фаз (1..8).
        \item Фаза 4: \texttt{\#include}, \texttt{\#if}, макро-экспансия, удаление неактивных веток.
        \item Фаза 7: построение AST, типизация, overload resolution, инстанцирование шаблонов.
        \item Следствие: на фазе 4 есть власть над текстом, но нет доступа к семантике C++.
    \end{itemize}
    \vspace{0.4em}
    Источники: \cite{cpp_draft_lex_phases,cpp_draft_cpp_pre}.
\end{frame}

\begin{frame}[fragile]{Инструментарий фазы 4: что именно умеет препроцессор}
\begin{lstlisting}
#define PI 3.14
#define MAX(a, b) ((a) > (b) ? (a) : (b))
#define LOG(fmt, ...) log(fmt, __VA_ARGS__)
#define STR(x) #x
#define CAT(a, b) a##b
\end{lstlisting}
    \begin{itemize}
        \item \texttt{\#define}: подстановка токенов по шаблону.
        \item \texttt{\#}/\texttt{\#\#}: превращение токенов в строку и склейка токенов.
        \item \texttt{\#if/\#ifdef/\#elif}: ранний срез кода до запуска парсера C++.
    \end{itemize}
\end{frame}

\begin{frame}[fragile]{Variadic-макросы: висячая запятая и \texttt{\_\_VA\_OPT\_\_}}
\begin{lstlisting}
#define LOG_BAD(fmt, ...) std::printf(fmt, __VA_ARGS__)
LOG_BAD("Hello");
\end{lstlisting}
\begin{lstlisting}
#define LOG_OK(fmt, ...) std::printf(fmt __VA_OPT__(,) __VA_ARGS__)
LOG_OK("Hello");
LOG_OK("x=%d", 42);
\end{lstlisting}
    \begin{itemize}
        \item Исторический костыль: \texttt{, \#\#\_\_VA\_ARGS\_\_} (vendor extension).
        \item C++20 стандартизировал корректный механизм: \texttt{\_\_VA\_OPT\_\_}.
    \end{itemize}
    \vspace{0.3em}
    Источник: \cite{wg21_p0306r4}.
\end{frame}

\begin{frame}[fragile]{Условная компиляция и \texttt{\_\_has\_include}}
\begin{lstlisting}
#if __has_include(<optional>)
  #include <optional>
#else
  #include "fallback/optional.hpp"
#endif
\end{lstlisting}
    \begin{itemize}
        \item \texttt{\#if/\#ifdef} физически удаляют ветви до построения AST.
        \item \texttt{\_\_has\_include} (C++17) дал переносимую проверку наличия заголовка.
        \item Это базовый паттерн graceful degradation для библиотек.
    \end{itemize}
    \vspace{0.3em}
    Источники: \cite{cpp_draft_cpp_cond,wg21_p0061r1}.
\end{frame}

\begin{frame}[fragile]{\texttt{\#pragma} и \texttt{\#embed}: backdoor и эволюция}
\begin{lstlisting}
#if defined(__clang__)
  #pragma clang diagnostic push
  #pragma clang diagnostic ignored "-Wdeprecated-declarations"
#endif
legacy_api();
#if defined(__clang__)
  #pragma clang diagnostic pop
#endif
\end{lstlisting}
\end{frame}

\begin{frame}[fragile]{\texttt{\#pragma} и \texttt{\#embed}: backdoor и эволюция}
    \begin{itemize}
        \item \texttt{\#pragma}: прямой интерфейс к конкретному компилятору, использовать локально.
        \item \texttt{\#embed} (C++26): встраивание ресурсов без генерации гигантских \texttt{.cpp}-массивов.
        \item Итог: инструменты становятся умнее, но природа фазы 4 остается текстовой.
    \end{itemize}
    \vspace{0.3em}
    Источник: \cite{wg21_p1967r14}.
\end{frame}

\section{Часть 2: Хорошие паттерны}

\begin{frame}[fragile]{Ленивый control flow: логирование без лишней работы}
\begin{lstlisting}
#define LOG_DEBUG(logger, ...)                                   \
    do {                                                          \
        if ((logger).is_debug()) {                               \
            (logger).log(__VA_ARGS__);                           \
        }                                                         \
    } while(false)
\end{lstlisting}
    \begin{itemize}
        \item Функция вычисляет аргументы eagerly, даже если логгер их потом отбросит.
        \item Макрос переносит вычисление внутрь \texttt{if} и устраняет лишний runtime.
        \item \texttt{do\{...\}while(false)} делает макрос безопасным statement.
    \end{itemize}
\end{frame}

\begin{frame}[fragile]{X-Macro: один источник правды для enum + to\_string}
\begin{lstlisting}
#define FOR_EACH_STATE(X) \
    X(IDLE)               \
    X(RUN)                \
    X(STOP)

enum class State { #define X(s) s, FOR_EACH_STATE(X) #undef X };
\end{lstlisting}
\end{frame}

\begin{frame}[fragile]{X-Macro: один источник правды для enum + to\_string}
\begin{lstlisting}
constexpr const char* to_string(State s) {
    switch (s) { #define X(s) case State::s: return #s; FOR_EACH_STATE(X) #undef X }
    return "UNKNOWN";
}
\end{lstlisting}
    \begin{itemize}
        \item Генерация повторяющихся деклараций из одного списка.
        \item Цена: низкая читаемость, поэтому нужен строгий стиль владения макросом.
    \end{itemize}
\end{frame}


\begin{frame}[fragile]{Платформенный слой: где \texttt{if constexpr} бессилен}
\begin{lstlisting}
#if defined(_WIN32)
  #include <windows.h>
#elif defined(__linux__)
  #include <sys/epoll.h>
#else
  #error "Unsupported platform"
#endif
\end{lstlisting}
    \begin{itemize}
        \item Конфликтующие платформенные API часто нельзя даже парсить вместе в одном TU.
        \item \texttt{if constexpr} работает на фазе 7, а значит слишком поздно для этой задачи.
        \item \texttt{\#if}-срез до парсинга здесь является единственным практичным механизмом.
    \end{itemize}
\end{frame}

\section{Часть 3: Цена власти}

\begin{frame}{Где начинается саботаж проекта}
    \begin{itemize}
        \item Проблема не в самом препроцессоре, а в неправильной зоне применения.
        \item Опасно использовать макросы для задач бизнес-логики, типов и layout-контрактов.
        \item Три системные границы: отсутствие гигиены, copy-paste модель, деградация диагностики.
    \end{itemize}
\end{frame}

\begin{frame}[fragile]{Предел 1: нет гигиены и границ}
\begin{lstlisting}
// packet.hpp
struct Packet {
    std::uint32_t id;
#ifdef ENABLE_PROFILING
    std::uint64_t ts;
#endif
    std::uint32_t payload;
};
// sender.cpp:   -DENABLE_PROFILING
// receiver.cpp:
\end{lstlisting}
\end{frame}

\begin{frame}[fragile]{Предел 1: нет гигиены и границ}
    \begin{itemize}
        \item Макросы протекают через include-граф и игнорируют namespace-границы.
        \item Результат: поздние и дорогие ODR-баги на линковке и в рантайме.
    \end{itemize}
\end{frame}

\begin{frame}{Предел 2: The Copy-Paste Curse}
    \begin{itemize}
        \item Фаза 4 не переиспользует логику как функции/шаблоны, а размножает токены.
        \item Следствие 1: Macro Expansion Hell в диагностике.
        \item Следствие 2: скрытый \texttt{.rodata} bloat из-за строк \texttt{\_\_FILE\_\_}, \texttt{\_\_PRETTY\_FUNCTION\_\_}, \texttt{\#expr}.
        \item Чем глубже иерархия макросов, тем хуже локальность ошибки и сложнее отладка.
    \end{itemize}
\end{frame}

\begin{frame}[fragile]{Хороший юзкейс, который легко сломать}
\begin{lstlisting}
#define SMART_ASSERT(expr)                                      \
    do {                                                        \
        if (!(expr)) {                                          \
            FatalLog(__FILE__, __LINE__, #expr, "false");      \
            std::abort();                                       \
        }                                                       \
    } while(0)
\end{lstlisting}
    \begin{itemize}
        \item Плюс: \texttt{\#expr} дает строку выражения, \texttt{\_\_FILE\_\_}/\texttt{\_\_LINE\_\_} дают контекст.
        \item Минус: каждый вызов порождает новые строковые литералы и новый expansion trace.
        \item На малом масштабе это удобно; на большом масштабе это источник системной стоимости.
    \end{itemize}
\end{frame}

\begin{frame}[fragile]{Эффект бабочки: Binary Bloat и Macro Expansion Hell}
    \begin{lstlisting}
macro_hell_demo.cpp:29:23: error: 'string' is not a member of 'std'
...
macro_hell_demo.cpp:25:28: note: in expansion of macro 'CHECK_16'
macro_hell_demo.cpp:29:5: note: in expansion of macro 'SMART_ASSERT'
note: 'std::string' is defined in header '<string>'
    \end{lstlisting}
    \begin{itemize}
        \item Ошибка появляется в фазе 7, но объясняется через шаги мутации текста фазы 4.
        \item Линкер склеивает только одинаковые строки; уникальные \texttt{\#expr} остаются в \texttt{.rodata}.
        \item Итог: хуже читаемость ошибок + рост бинарника без роста бизнес-ценности.
    \end{itemize}
\end{frame}

\begin{frame}[fragile]{Лекарство: перенос метаданных в фазу 7}
\begin{lstlisting}
#include <source_location>

inline void smart_assert(bool expr, const char* expr_str,
    std::source_location loc = std::source_location::current()) {
    if (!expr) {
        FatalLog(loc.file_name(), loc.line(), expr_str);
        std::abort();
    }
}
\end{lstlisting}
\end{frame}

\begin{frame}[fragile]{Лекарство: перенос метаданных в фазу 7}
    \begin{itemize}
        \item \texttt{\_\_FILE\_\_}/\texttt{\_\_LINE\_\_} уходят из макроса в типобезопасный API.
        \item Диагностика локализуется в функции, а не в этажах макро-экспансии.
        \item Для \texttt{\#expr} можно оставить тонкую макро-обертку, тяжелую логику держать в C++.
    \end{itemize}
\end{frame}

\begin{frame}{современномый C++: перенос логики на фазу 7}
    \begin{itemize}
        \item Вместо \texttt{\#define}-констант: \texttt{constexpr} и inline-переменные.
        \item Вместо тяжелых макросов для метаданных: \texttt{std::source\_location}.
        \item Вместо макро-ветвлений внутри логики: \texttt{if constexpr}, concepts, modules.
        \item Правило: фазе 4 оставлять только то, что невозможно решить после парсинга.
    \end{itemize}
\end{frame}

\section{Часть 4: Сравнение подходов (Rust и Go)}

\begin{frame}[fragile]{Rust: сохранить мощь, убрать хаос}
\begin{lstlisting}[language=Rust]
#[cfg(target_os = "linux")]
const TARGET: &str = "linux";
#[cfg(not(target_os = "linux"))]
const TARGET: &str = "not-linux";
\end{lstlisting}
\end{frame}

\begin{frame}[fragile]{Rust: сохранить мощь, убрать хаос}
\begin{lstlisting}[language=Rust]
macro_rules! define_states {
    ($($state:ident),+ $(,)?) => {
        enum State { $($state),+ }
        impl State {
            fn as_str(self) -> &'static str {
                match self { $(State::$state => stringify!($state)),+ }
            }
        }
    };
}
define_states!(Idle, Running, Failed);
\end{lstlisting}
\end{frame}

\begin{frame}[fragile]{Rust: сохранить мощь, убрать хаос}
    \begin{itemize}
        \item \texttt{macro\_rules!} сопоставляет и разворачивает токены через matcher/transcriber.
        \item В правилах доступен фрагмент \texttt{tt} (TokenTree), а гигиена описана как mixed-site.
        \item \texttt{\#[cfg]} и \texttt{cfg!} задают условную компиляцию в самом языке.
    \end{itemize}
    \vspace{0.3em}
    Источники: \cite{rust_ref_macros_by_example,rust_ref_conditional_compilation}.
\end{frame}

\begin{frame}[fragile]{Go: радикальный отказ от препроцессора}
\begin{lstlisting}[language=Go]
//go:build linux && amd64
package main

func platformName() string { return "linux-amd64" }
\end{lstlisting}
\begin{lstlisting}
//go:generate stringer -type=State
type State int
const (
    Idle State = iota
    Running
    Failed
)
\end{lstlisting}
\end{frame}

\begin{frame}[fragile]{Go: радикальный отказ от препроцессора}
\begin{itemize}
        \item Условная сборка задается build constraints (\texttt{//go:build}) на уровне файлов.
        \item Кодогенерация задается директивами \texttt{//go:generate}.
        \item \texttt{go generate} не является частью \texttt{go build}: это явный отдельный шаг.
    \end{itemize}
    \vspace{0.3em}
    Источники: \cite{go_cmd_build_constraints,go_cmd_generate,go_faq_headers}.
\end{frame}

\section{Заключение}

\begin{frame}{Сравнительная таблица подходов}
    \centering
    \small
    \begin{tabular}{p{2cm}p{2.5cm}p{2.5cm}p{2.5cm}}
        \textbf{Критерий} & \textbf{C++} & \textbf{Rust} & \textbf{Go} \\
        \hline
        \textbf{Фаза} & До парсинга & После парсинга & Вне компилятора \\
        \hline
        \textbf{Гигиена} & Нет & Есть & N/A \\
        \hline
        \textbf{Диагностика} & Expansion Hell & TokenTree & Прямые ошибки \\
    \end{tabular}
\end{frame}

\begin{frame}{Финальная формула доклада}
    \begin{itemize}
        \item Эволюция C++ не отменяет препроцессор, а сужает его зону ответственности.
        \item Фаза 4 сильна в раннем срезе кода, \texttt{\#pragma} и тонких zero-overhead обертках.
        \item Фаза 4 слаба в задачах типов, гигиены и диагностики: их лучше решать на фазе 7.
        \item Правило архитектора: минимально мощный механизм, достаточный для задачи.
    \end{itemize}
    \vspace{0.5em}
    \centering \textbf{Уважайте фазу 4, но не заставляйте ее играть роль фазы 7.}
\end{frame}

\section{Источники}

\begin{frame}[allowframebreaks]{Источники}
    \nocite{cpp_draft_lex_phases,cpp_draft_cpp_pre,cpp_draft_cpp_cond,cppref_preprocessor_replace,wg21_p0061r1,wg21_p0306r4,wg21_p2334r1,wg21_p2437r1,wg21_p1967r14,rust_ref_macros_by_example,rust_ref_conditional_compilation,go_cmd_build_constraints,go_cmd_generate,go_faq_headers}
    \bibliographystyle{unsrt}
    \bibliography{ref}
\end{frame}

\end{document}
